\begin{theorem}[Proved: the square-zero criterion with arbitrary dyadic ramification]
\label{C:ramified_squarezero_criterion}
Let $A$ be any DVR of characteristic zero and residue characteristic
two, let $\pi$ be a uniformizer, and let $J$ be any $A$-module.
For $k=A\ltimes J$, the three conditions of
Theorem~\ref{C:squarezero_intrinsic_criterion} remain equivalent:
the full polynomial induction-ring answer, its degree-two
surjectivity, and the bijectivity of $2:J\to J$.
There is no restriction on the ramification index or on $J$.
\end{theorem}
\begin{proof}
Write $2=v\pi^e$ with $v\in A^\times$ and $e\ge1$.
The sufficient direction and the determinant obstruction when
$J[2]\ne0$ are exactly those of
Theorem~\ref{C:squarezero_intrinsic_criterion}.
Suppose therefore that $2$ is injective on $J$ but not surjective.
Then $J$ is flat over $A$, and $N=J/2J$ is a nonzero flat module
over the artinian local ring $B=A/2A$.
We record directly why $N$ is free, allowing arbitrary rank.
Lift a basis of $N/\pi N$ to obtain a map $B^{(I)}\to N$.
Its cokernel $C$ satisfies $C=\pi C$, hence $C=\pi^eC=0$.
For its kernel $K$, flatness of $N$ makes
$K/\pi K\to B^{(I)}/\pi B^{(I)}$ injective; its image is zero
by the chosen basis, so $K=\pi K=\pi^eK=0$.
Thus a basis coefficient gives a $B$-linear functional
$\ell:J/2J\to B$ and an $x\in J$ with $\ell(x)=1$.

For a finite free $A$-lattice $L$ with involution, its two Tate
groups are finite-length $B$-modules. Flatness gives the natural
comparison
\[
 \widehat H^i(C_2,L\otimes_AJ)
    \simeq \widehat H^i(C_2,L)\otimes_B(J/2J).
\]
For a representation over $k$, write $T=T_0+U$ as in
Lemma~\ref{C:linear_nil_Bockstein}. The same coefficient identities
and the functional $\ell$ give natural odd maps $b_\ell$ on these
Tate groups, with the same anticommutation with connecting maps.
Put $u=b_\ell^2$.

Let $F=A/\pi A$. Finite-length $A[u]$-modules decompose exactly
into their primary parts at the maximal ideals $(\pi,q(u))$,
where $q\in F[u]$ is monic irreducible. To see the assertion
without choosing lifts of roots, use the finite support of a
finite-length module and the Chinese remainder idempotents for
sufficient powers of its finitely many distinct maximal ideals.
These idempotents are compatible with every module map, so this
decomposition preserves exact sequences. For $q\ne u$, the
operator $u$ is invertible on the indicated primary part.
The exact Tate six-cycle therefore gives the same exact
three-cycle as in Lemma~\ref{C:epsilon_Bockstein}. The length
of the even primary part, now its composition length as an
$A[u]$-module, is additive modulo two. Denote this invariant
by $\chi_{q,\ell}$.

Put $c=2/\pi\in A$ and consider the integral matrices
\[
 T_0=\begin{pmatrix}1&\pi\\0&-1\end{pmatrix},\qquad
 U=\begin{pmatrix}1&0\\-c&-1\end{pmatrix}.
\]
They satisfy $T_0^2=U^2=1$ and $T_0U+UT_0=0$.
Thus $T=T_0+xU$ is an involution over $k$.
The plus cycles of $T_0$ are generated by $e_1$, and the
minus cycles by $w=e_1-ce_2$. The norm image on the plus
cycles is $(2,\pi)Ae_1=\pi Ae_1$, and the difference image
on the minus cycles is $\pi Aw$. Consequently both Tate
groups are $F$, including when $e=1$.
Since $Ue_1=w$ and $Uw=e_1$, the two maps $b_\ell$ are
identities. Hence
\[
 \chi_{u-1,\ell}([T]-1-s)=1.
\]
The invariant vanishes on both constant lines. This proves failure
of the degree-two polynomial answer in every remaining case.
The class is generically zero because the two eigenspaces become
free rank-one modules over the local rational square-zero ring.
It has exact order two: $T_0U$ is invertible and anticommutes
with $T_0$ and $U$, so it identifies $T$ with $-T$; the
evaluation sequence then gives twice the class equal to zero.
\end{proof}
